\secciongrande{Capacidades y limitaciones de ejecución: ¿por qué JavaScript en el cliente}

\resumen{
  Aunque la gestión dinámica principal ocurre en el servidor, JavaScript es esencial en el cliente para mejorar la usabilidad. Sus capacidades incluyen la preparación de datos, la navegación sin dependencia total de JS, y la actualización dinámica del contenido recibido del servidor.
}

\subtituloA{Razones para usar JavaScript en el cliente}

\lista{
  \textbf{Preparación de datos:} JavaScript no puede escribir archivos en el servidor; su función es preparar datos para su envío o cederlos a otro lenguaje del servidor. |
  \textbf{Navegación resiliente:} El sitio debe ser navegable sin JavaScript; este se usa para modernizar la experiencia, pero no debe ser un requisito absoluto para la funcionalidad básica. |
  \textbf{Actualización dinámica:} Permite examinar datos recibidos del servidor y actualizar selectivamente secciones de la página web sin recargarla completamente.
}
